\section{The model potential of a twisted canonical ratio line}

Let $\mathcal X=(\mathcal X\to\mathbb A^1,\ \mathbb G_m\text{-action},\
\mathcal X_1\cong X)$ be a test configuration of a proper $k$-scheme $X$, let
$U=\mathcal X|_{t\neq0}$ be its noncentral open, let $\pi:\mathcal X\to X$ be
a morphism with $\pi|_U$ equal to the projection of the trivialisation
$U\cong X\times\mathbb G_m$ and $\pi|_{\mathcal X_1}$ equal to the fibre-one
identification, and let $L$ be invertible on $X$.  Let $K,J$ be invertible
ideal sheaves on $\mathcal X$, invariant under the action, whose supports miss
$U$.  The \emph{twisted ratio line} is
$M=\pi^*L\otimes K\otimes J^\vee$, and its \emph{canonical ratio section}
along a morphism $f$ factoring through $U$ is
$\sigma_f=(f^*K\to\mathcal O)^{-1}\otimes f^*(J\hookrightarrow\mathcal O)^\vee$.
We compute the native model potential (Definition~\ref{def:native-tc-original-model-potential})
of $M$.

\begin{lemma}[Calculus of pulled and tensor sections]
\label{lem:twisted-ratio-section-calculus}
\lean{EnvelopeProperty.PLOriginalModelCompatibility.TwistedRatioCalculus.pullbackSectionHom_map, EnvelopeProperty.PLOriginalModelCompatibility.TwistedRatioCalculus.pullbackSectionHom_comp_inv, EnvelopeProperty.PLOriginalModelCompatibility.TwistedRatioCalculus.pullbackSectionHom_congr, EnvelopeProperty.PLOriginalModelCompatibility.TwistedRatioCalculus.pullbackSectionHom_scalar, EnvelopeProperty.PLOriginalModelCompatibility.TwistedRatioCalculus.pullbackSectionHom_unit_comp, EnvelopeProperty.PLOriginalModelCompatibility.TwistedRatioCalculus.tensorSectionHom_rightUnit, EnvelopeProperty.PLOriginalModelCompatibility.TwistedRatioCalculus.pulledTensorSection, EnvelopeProperty.PLOriginalModelCompatibility.TwistedRatioCalculus.pulledTensorSection_comp_hom, EnvelopeProperty.PLOriginalModelCompatibility.TwistedRatioCalculus.pullbackSectionHom_pulledTensorSection, EnvelopeProperty.PLOriginalModelCompatibility.TwistedRatioCalculus.pulledTensorSection_congr, EnvelopeProperty.PLOriginalModelCompatibility.TwistedRatioCalculus.pulledTensorSection_pullbackIsoAlong}
\leanok
For a morphism $f$ and sections $a$ of $f^*T$, $b$ of $f^*R$ write
$\langle a,b\rangle_f$ for the section of $f^*(T\otimes R)$ obtained from
$a\otimes b$ through the tensorator of $f^*$ and the comparison of the
monoidal and explicit tensor products.  Then:
(i) $f^*$ of sections is functorial and compatible with the composition and
equality comparisons of pullbacks, and a scalar factor pulls back to the
pulled scalar;
(ii) $g^*\langle a,b\rangle_f$, transported by $g^*f^*\cong(fg)^*$, is
$\langle g^*a,g^*b\rangle_{fg}$ transported factorwise;
(iii) if a comparison $\alpha:f^*(T\otimes R)\cong g^*(T\otimes R)$ has the
tensor law $\alpha\langle\cdot,\cdot\rangle_f=\langle\alpha_T\cdot,\alpha_R\cdot\rangle_g$,
then its base change along any $h$ sends $\langle a,b\rangle_{hf}$ to
$\langle h^*\alpha_T\, a,h^*\alpha_R\, b\rangle_{hg}$;
(iv) a tensor section with a scalar second factor contracts to the scalar
multiple of the first factor.
\end{lemma}
\begin{proof}
(i) A pulled section is the inverse structure-sheaf comparison followed by
the pulled morphism; functoriality is functoriality of $f^*$, the
composition rule is the coherence of the structure-sheaf comparison with
$g^*f^*\cong(fg)^*$, and equality comparisons are identities after
substituting the equality.  (ii) The pulled pure tensor section is the pure
tensor of the pulled sections followed by the tensorator; the composition
comparison is a monoidal natural isomorphism, so it carries the iterated
tensorator to the tensorator of the composite, and it is natural for the
explicit tensor comparison.  (iii) Write $h^*\alpha$ between the composition
comparisons and apply (ii) on both sides; the tensor law of $\alpha$ and
naturality of the tensorator of $h^*$ move $\alpha$ onto the factors.
(iv) Every endomorphism of the structure sheaf is multiplication by a global
function; scalars pass through tensor sections, and the tensor section with
the unit section followed by the right unitor is the first factor.
\end{proof}

\begin{lemma}[Naturality of the canonical ratio section]
\label{lem:twisted-ratio-canonical-section-natural}
\lean{EnvelopeProperty.PLOriginalModelCompatibility.TwistedRatioCalculus.canonicalPulledRatioSection_eq_pulledTensorSection, EnvelopeProperty.PLOriginalModelCompatibility.TwistedRatioCalculus.pullbackSectionHom_canonicalPulledRatioSection, EnvelopeProperty.PLOriginalModelCompatibility.TwistedRatioCalculus.pullbackSectionHom_twistedCanonicalSection, EnvelopeProperty.PLOriginalModelCompatibility.TwistedRatioCalculus.twistedCanonicalSection_comp_inv, EnvelopeProperty.PLOriginalModelCompatibility.TwistedRatioCalculus.twistedCanonicalSection_congr, EnvelopeProperty.PLOriginalModelCompatibility.TwistedRatioCalculus.pullbackResidueSection_twistedCanonical, EnvelopeProperty.PLOriginalModelCompatibility.TwistedRatioCalculus.transportedResidueSection_twistedCanonical}
\uses{lem:twisted-ratio-section-calculus,thm:original-pl-canonical-ratio-section-transport}
\leanok
If $K$ pulls back to invertible inclusions along $g$ and $fg$, then
$f^*\sigma_g=\sigma_{fg}$ and $f^*\langle g^*s,\sigma_g\rangle_g=\langle (fg)^*s,\sigma_{fg}\rangle_{fg}$.
In particular the residue pullback of the fibre of
$\langle s,\sigma\rangle$ at $g(y)$ is the fibre at $y$ of
$\langle g^*s,\sigma_g\rangle_g$, and its transport by any comparison
$\Phi:g^*M\cong P$ is the fibre of $\Phi\langle g^*s,\sigma_g\rangle_g$.
\end{lemma}
\begin{proof}
The canonical ratio section is $\langle (f^*K\to\mathcal O)^{-1},
f^*\delta\rangle_f$ with $\delta$ the transpose of $J\hookrightarrow\mathcal O$.
The pulled numerator inverse is the numerator inverse along the composite,
because pulled inclusions compose with the structure comparison; the
denominator part is a pulled section.  Apply Lemma~\ref{lem:twisted-ratio-section-calculus}(ii).
The residue statements follow by writing the residue comparison as a
composition comparison followed by an equality comparison.
\end{proof}

\begin{lemma}[Structural laws of the tensor-form comparisons]
\label{lem:twisted-ratio-comparison-laws}
\lean{EnvelopeProperty.PLOriginalModelCompatibility.TwistedRatioCalculus.canonicalTensorPullbackIso_comparison_tensor_law, EnvelopeProperty.PLOriginalModelCompatibility.TwistedRatioCalculus.invariantComparison_pullbackSectionHom, EnvelopeProperty.PLOriginalModelCompatibility.TwistedRatioCalculus.pulledTensorSection_unitTrivialization, EnvelopeProperty.NativeTestConfiguration.TestConfiguration.twistedRatio_twistLaw_of_iso, EnvelopeProperty.NativeTestConfiguration.TestConfiguration.twistedRatio_fibreOneLaw_of_iso}
\uses{lem:twisted-ratio-section-calculus}
\leanok
A comparison of explicit tensor products built factorwise from $\alpha_T$ and
$\alpha_R$ through the canonical tensor-pullback isomorphisms has the tensor
law of Lemma~\ref{lem:twisted-ratio-section-calculus}(iii); the comparison of
a pullback along an invariant projection fixes pulled sections; and a
factorwise fibre-one identification with trivialised second factor $e$ sends
$\langle a,b\rangle$ to $(e\circ b)\cdot a$.  The last two laws hold for every
comparison equal, as an isomorphism, to the invariant-projection comparison,
respectively to the factorwise fibre-one identification.
\end{lemma}
\begin{proof}
Unfold the canonical tensor-pullback isomorphisms; the explicit and monoidal
tensor comparisons are natural, so the comparison is the tensorator,
$\alpha_T\otimes\alpha_R$, and the inverse tensorator.  The invariant
comparison is a composite of composition and equality comparisons, which
fix pulled sections by Lemma~\ref{lem:twisted-ratio-section-calculus}(i).
The last law is Lemma~\ref{lem:twisted-ratio-section-calculus}(iv) after the
right-unit identification of the explicit tensor with the unit.
\end{proof}

\begin{proposition}[Transport of the twisted canonical section]
\label{prop:twisted-ratio-section-transport}
\lean{EnvelopeProperty.PLOriginalModelCompatibility.TwistedRatioCalculus.comp_trans_seven_hom, EnvelopeProperty.NativeTestConfiguration.TestConfiguration.unitOpenLineIso_eq_explicit, EnvelopeProperty.NativeTestConfiguration.TestConfiguration.twistedRatio_inclusion_isIso, EnvelopeProperty.NativeTestConfiguration.TestConfiguration.twistedRatio_section_step0, EnvelopeProperty.NativeTestConfiguration.TestConfiguration.twistedRatio_section_step1, EnvelopeProperty.NativeTestConfiguration.TestConfiguration.twistedRatio_section_step2, EnvelopeProperty.NativeTestConfiguration.TestConfiguration.twistedRatio_section_step3, EnvelopeProperty.NativeTestConfiguration.TestConfiguration.twistedRatio_section_step4, EnvelopeProperty.NativeTestConfiguration.TestConfiguration.twistedRatio_section_step5, EnvelopeProperty.NativeTestConfiguration.TestConfiguration.twistedRatio_section_step6, EnvelopeProperty.NativeTestConfiguration.TestConfiguration.twistedRatio_unitOpenLineIso_section, EnvelopeProperty.NativeTestConfiguration.TestConfiguration.twistedRatio_productTotalLineIso_section}
\uses{lem:twisted-ratio-canonical-section-natural,lem:twisted-ratio-comparison-laws,cor:original-pl-group-canonical-ratio-section}
\leanok
Let the action comparison of $M$ have the tensor law with factors
$\alpha_T$ (fixing pulled sections of $\pi^*L$) and $\alpha_R$ (fixing the
canonical ratio section along the action lift of $U$), and let the fibre-one
identification of $M$ be factorwise with trivialisation $e_R$ of the ratio on
$\mathcal X_1$.  For every global section $s$ of $L$, the whole noncentral
line comparison sends $\langle \iota_U^*\pi^*s,\sigma_{\iota_U}\rangle$ to
$\varepsilon\cdot w^*s$, where $\varepsilon$ is the pulled value of
$e_R\circ\sigma_{\mathcal X_1}$; transported to $X\times\mathbb G_m$, the
corresponding section of the pulled $M$ is sent to $\mathrm{pr}_1^*(\varepsilon_X s)$,
with $\varepsilon_X$ a global unit of $X$.
\end{proposition}
\begin{proof}
The noncentral comparison is the equality comparison $\iota_U=\ell\circ a$
($\ell$ the action lift, $a$ the action), the base change of the action
comparison along $\ell$, the equality comparison $\ell\circ p=\zeta\circ
\iota_1$ ($p$ the projection, $\zeta$ the fibre projection), the inverse
composition comparison, $\zeta^*$ of the fibre-one identification, and the
composition and equality comparisons to $w=\pi\circ\iota_U$.  Follow the
section through: the equality comparisons and
Lemma~\ref{lem:twisted-ratio-canonical-section-natural} keep its form, the
action step is Lemma~\ref{lem:twisted-ratio-section-calculus}(iii) with the
invariance of the canonical section
(Corollary~\ref{cor:original-pl-group-canonical-ratio-section}) and of pulled
sections, the fibre-one step is the last law of
Lemma~\ref{lem:twisted-ratio-comparison-laws}, and the remaining comparisons
compose pulled sections.  On $X\times\mathbb G_m$ the comparison is the
pullback along $U\cong X\times\mathbb G_m$; the scalar pulls back along
$X\times\mathbb G_m\to U\to\mathcal X_1\cong X$, which is $\mathrm{pr}_1$
followed by the inverse fibre-one identification.  The unit $\varepsilon_X$
is invertible because the canonical ratio section is a frame off both
supports (Theorem~\ref{thm:original-canonical-ratio-section-frame}).
\end{proof}

\begin{theorem}[The model potential of a twisted ratio line]
\label{thm:twisted-ratio-model-potential}
\lean{EnvelopeProperty.NativeTestConfiguration.TestConfiguration.twistedRatio_mapOverBase_productGauss, EnvelopeProperty.NativeTestConfiguration.TestConfiguration.gaussCompatibleNormFactor_twistedRatio, EnvelopeProperty.NativeTestConfiguration.TestConfiguration.gaussNativeNormFactor_twistedRatio, EnvelopeProperty.NativeTestConfiguration.TestConfiguration.gaussNativeModelPotential_twistedRatio}
\uses{prop:twisted-ratio-section-transport,thm:pl-original-native-factor-arbitrary-residue-vector,def:native-tc-original-model-potential,thm:original-canonical-ratio-section-frame}
\leanok
Under the hypotheses of Proposition~\ref{prop:twisted-ratio-section-transport},
if $\pi\colon\mathcal X\to X$ lies over $k\to k[t]$, then for every
$0<r\le1$ and every trivially valued point $x$ of $X$ with a global section $s$
of $L$ such that $\|s\|(x)\neq0$,
\[
  \phi_{\mathcal X,M}(x,r)=-\log\|\sigma\|(z_{x,r}),
\]
where $z_{x,r}$ is the radius-$r$ Gauss point of $\mathcal X$ over $x$ and
$\sigma$ is the canonical ratio section at its residue field.
\end{theorem}
\begin{proof}
Let $c$ be the defining factor of the potential for any centre-open frame of
$L$ at $x$.  By the residue formula of the factor
(Theorem~\ref{thm:pl-original-native-factor-arbitrary-residue-vector}), every
vector $v$ of the fibre of $M$ at $z_{x,r}$ satisfies
$\|v\|=c\,|\mathrm{coef}(\Phi v)|$, the coefficient being taken in the
reference frame at the product Gauss point.  Take $v$ the fibre of
$\langle\pi^*s,\sigma\rangle$.  Its model norm is
$\|s\|(x)\,\|\sigma\|(z_{x,r})$: model norms are multiplicative on tensor
sections and the model norm of a pulled section is the model norm at the
image point, which is $x$.  By
Proposition~\ref{prop:twisted-ratio-section-transport} and
Lemma~\ref{lem:twisted-ratio-canonical-section-natural}, $\Phi v$ is the fibre
of $\mathrm{pr}_1^*(\varepsilon_X s)$; its only Laurent coefficient is
$\varepsilon_X s$ in degree $0$, so its reference coefficient has value
$\|\varepsilon_Xs\|(x)=\|s\|(x)$, a unit of the structure sheaf having model
norm one.  Hence $\|s\|(x)\|\sigma\|(z_{x,r})=c\,\|s\|(x)$ and
$c=\|\sigma\|(z_{x,r})$.  Taking $-\log$ gives the potential.
\end{proof}
